% 2106 Proof Template
% Last revised Spring 2019

\documentclass[12pt]{amsart}
\usepackage{amsmath, amssymb, amsthm, amsfonts}
\usepackage{dsfont, bbm, marvosym}
\pagestyle{empty}

% ***** PREAMBLE *****
% Do not modify anything in this part.  In particular, DO NOT modify the margins.

% These pre-defined commands give you easy access to "blackboard bold" letters.  Feel free to add more!
\newcommand{\R}{\ensuremath{\mathbb{R}}}
\newcommand{\Z}{\ensuremath{\mathbb{Z}}}
\newcommand{\Q}{\ensuremath{\mathbb{Q}}}
\newcommand{\N}{\ensuremath{\mathbb{N}}}

\renewcommand{\qedsymbol}{\rule{0.6em}{0.6em}}

% This command defines the "solution" environment that you'll use for your proofs.  Please do NOT modify this.
\newenvironment{solution}[2]{\noindent \textbf{{#1}}. {#2} \\  \begin{proof}}{\end{proof}}

% ***** IDENTIFYING INFORMATION *****

\title{Second Partial Portfolio- Math 250- Spring 2026}

\author{Leo Winston}

\date{April 22, 2026}

% ***** MAIN DOCUMENT ***** 

\begin{document}

\maketitle
\section{Set Theoretic Proof}
{\begin{center}(Homework 3 $\#14$)\end{center}}
\textbf{Theorem:}{ Let $A, B$, and $C$ be arbitrary sets. If $A-C \nsubseteq A-B$, then $B \nsubseteq C$.}
\begin{proof}
Let $A, B, C$ be arbitrary sets. We will proceed by proving the contrapositive. That is, we will show $B \subseteq C$ implies $A-C \subseteq A-B$. Assume $B \subseteq C$, and let $x$ be an arbitrary element of $A - C$. By definition of set difference, $x \in A - C$ implies $x \in A$ and $x \not \in C$. By definition of a subset, the assumption $B \subseteq C$ means every element of $B$ is also an element of $C$. Consequently, since $x \notin C$, it logically follows $x \notin B$.
Now we have established that $x \in A$ and $x \notin B$. By definition of set difference, this implies $x \in A - B$. So, by definition of subset, $A - C \subseteq A - B$. Thus, we have proven the contrapositive. Therefore, the original statement holds: if $A - C \not\subseteq A - B$, then $ B \not\subseteq C$.
\end{proof}
\newpage

\section{Proof by Induction}
\begin{center}(Exam 2 Review $\#11$)\end{center}

\textbf{Theorem:} If $n \in \mathbb{N}$, then 
$1^3+2^3+3^3+4^3+\dots+n^3=\dfrac{n^2(n+1)^2}{4}$.

\begin{proof}
Let $n \in \mathbb{N}$, and let $P(n)$ be the statement
\[
    \sum_{i=1}^{n} i^3 = 1^3 + 2^3 + \cdots + n^3 = \frac{n^2(n+1)^2}{4}.
\]

\textbf{Base Case:} Let $n = 1$. Then,
\[
    \sum_{i=1}^{1} i^3 = 1^3 = 1,
\]
and
\[
    \frac{n^2(n+1)^2}{4} = \frac{1^2(1+1)^2}{4} = \frac{1 \cdot 4}{4} = 1.
\]
Since both sides equal $1$, $P(1)$ holds.

\textbf{Inductive Hypothesis:} Suppose $P(k)$ holds for some $k \in \mathbb{N}$.
That is, suppose
\[
    \sum_{i=1}^{k} i^3 = \frac{k^2(k+1)^2}{4}.
\]

\textbf{Inductive Step:} We want to show that $P(k+1)$ holds. By evaluating the sum up to $k+1$ and separating the final term, we may substitute our inductive hypothesis:
\begin{align*}
    \sum_{i=1}^{k+1}i^3 
        &= \left(\sum_{i=1}^{k}i^3\right) + (k+1)^3 \\
        &= \frac{k^2(k+1)^2}{4} + (k+1)^3 \quad \text{(by the inductive hypothesis)} \\
        &= (k+1)^2 \left(\frac{k^2}{4} + (k+1)\right) \\
        &= (k+1)^2 \left(\frac{k^2 + 4(k+1)}{4}\right) \\
        &= (k+1)^2 \left(\frac{k^2 + 4k + 4}{4}\right).
\end{align*}
Finally, by factoring $k^2+4k+4 = (k+2)^2$, we obtain:
\[
    \sum_{i=1}^{k+1}i^3 = \frac{(k+1)^2(k+2)^2}{4}.
\]

\nopagebreak[4]
This is precisely the statement $P(k+1)$. Since $P(1)$ holds, and $P(k)$ implies $P(k+1)$ for an arbitrary $k \in \mathbb{N}$, it follows by the principle of mathematical induction that $P(n)$ holds for all $n \in \mathbb{N}$. That is,
\[
    \sum_{i=1}^{n} i^3 = \frac{n^2(n+1)^2}{4}. \qedhere
\]
\end{proof}
\newpage
\section{Equivalence Relation Proof}
\begin{center}(Homework 8 $\#5$)\end{center}

\textbf{Theorem:} Let $R$ be the relation on $\mathbb{Z}$ defined by $aRb$ if $a^2-b^2$ is divisible by $4$. The relation $R$ is an equivalence relation on $\mathbb{Z}$.

\begin{proof} 
Let $R$ be the relation on $\mathbb{Z}$ defined by $aRb$ if $4 \mid (a^2-b^2)$. We will prove that $R$ is an equivalence relation by showing it is reflexive, symmetric, and transitive. 

\vspace{0.5cm}
\noindent \textbf{Reflexive:} Let $a \in \mathbb{Z}.$ Since $a$ is an integer, $a^2$ is also an integer. Because any integer subtracted from itself is zero, we know $a^2 - a^2 = 0$. Note that $0=n(0)$ for any integer $n$. Thus $0 = 4(0)$, and since $0 \in \mathbb{Z}$, it follows by definition of divides that $4 \mid (a^2-a^2)$. Thus, $aRa$ by definition of $R$, so the relation $R$ is reflexive.

\vspace{0.5cm}
\noindent \textbf{Symmetric:} Let $a,b \in \mathbb{Z}$. Suppose $aRb$. By definition of $R$, $4 \mid (a^2-b^2)$. By definition of divides, $a^2-b^2=4k$ for some integer $k$. Multiplying both sides of the equation by $-1$ yields $b^2-a^2 = 4(-k)$. Since $k$ is an integer, $-k$ is also an integer. Thus, by definition of divides, $4 \mid (b^2 - a^2)$. Therefore, $bRa$ by definition of $R$, proving that $R$ is symmetric.

\vspace{0.5cm}
\noindent \textbf{Transitive:} Let $a,b,c \in \mathbb{Z}$. Suppose $aRb$ and $bRc$. By definition of $R$, there exist integers $k$ and $n$ such that $a^2-b^2=4k$ and $b^2-c^2=4n$. Adding these two equations yields:
$$(a^2-b^2) + (b^2-c^2) = 4k + 4n$$
By simplifying the left-hand side and factoring the right-hand side, we obtain $a^2 - c^2 = 4(k+n)$. Since $k$ and $n$ are integers, their sum $k+n$ is also an integer due to closure under addition. Thus, by definition of divides, $4 \mid (a^2-c^2)$. Therefore, $aRc$ by definition of $R$, proving that $R$ is transitive.

Because the relation $R$ is reflexive, symmetric, and transitive, it is an equivalence relation on $\mathbb{Z}$.
\end{proof}
\newpage
\section{Injectivity and Surjectivity Proof}
\begin{center}(Homework 10 $\#1.e$)\end{center}
\textbf{Theorem:} The function $f: \mathbb{R} \rightarrow \mathbb{R}$ defined by $f(x) = x^4+x^2$ is neither injective, nor surjective.

\begin{proof}
To show that $f$ is neither injective nor surjective, we provide counterexamples for both properties.

\vspace{0.5cm}
\noindent \textbf{$f$ is not injective:} 
Consider the elements $a=1$ and $b=-1$ in the domain of the function $f$:
\begin{align*}
f(1)&=(1)^4+(1)^2 \\
&=1+1 \\
&=2 \\
f(-1)&=(-1)^4+(-1)^2 \\
&=(-1)^2(-1)^2+(-1)(-1) \\
&=(1)(1) + 1 \\
&= 2
\end{align*}
$f(1) = f(-1)=2$, but $1 \ne -1$. By definition of function injectivity, since we have found two distinct elements in the domain that are mapped to the same element in the codomain, $f$ is not injective.

\vspace{0.5cm}
\noindent \textbf{$f$ is not surjective:} 
Consider the element $y = -1$ in the codomain of $f$. We observe that for any $x \in \mathbb{R}$, $x^4 \ge 0$ and $x^2 \ge 0$. It follows that:
\[
    f(x) = x^4 + x^2 \ge 0
\]
for all $x \in \mathbb{R}$. Since $f(x)$ is always non-negative, there is no $x \in \mathbb{R}$ such that $f(x) = -1$. Thus, $f$ is not surjective.

\vspace{0.5cm}
We have shown that if $a,b \in \mathbb{R}$ such that $f(a)=f(b)$, it does not imply $a = b$. We have also shown that $y \in \mathbb{R}$ does not imply that there is an $x \in \mathbb{R}$ such that $y = f(x)$. The function $f(x)=x^4+x^2$ is therefore neither injective nor surjective.
\end{proof}
\newpage
\section{Continuity Proof}
\begin{center}(Homework 10 $\#2$)\end{center}

\textbf{Theorem:} The function $f: \mathbb{R} \rightarrow \mathbb{R}$ defined by $f(x) = 2x^2+1$ is continuous at $x=2$.

\begin{proof}
Let $\epsilon > 0$. We choose $\delta = \min\left(1, \frac{\epsilon}{10}\right)$. Suppose $x \in \mathbb{R}$ such that $|x-2| < \delta$. We will show that $|f(x) - f(2)| < \epsilon$.

Because we chose $\delta \le 1$, we know that $|x-2| < \delta \le 1$, and thus $|x-2|$ must be strictly less than $1$. This implies:
\begin{align*}
    -1 < x - 2 &< 1 \\
    -1+4 < (x - 2)+4 &< 1 +4\\
    3 < x + 2 &< 5 \\
\end{align*}
Since $x+2$ is strictly between $3$ and $5$, it is positive, and therefore we can conclude that $|x+2| < 5$.

Now, we must evaluate $|f(x)-f(2)|$. We are given that $f(x)=2x^2+1$, so by the function definition, and factoring, we obtain:
\begin{align*}
    |f(x) - f(2)| &= |(2x^2+1) - (2(2)^2+1)| \\
    &= |2x^2+1-9| \\
    &= |2x^2-8| \\
    &= 2|x^2-4| \\
    &= 2|x-2||x+2| 
\end{align*}
Because we have established that $|x+2| < 5$, and because our choice of $\delta$ guarantees that $|x-2| < \frac{\epsilon}{10}$, we can apply these bounds:
\begin{align*}
    2|x-2||x+2| &< 2 \left(\frac{\epsilon}{10}\right) (5) \\
    &= 10 \left(\frac{\epsilon}{10}\right) \\
    &= \epsilon
\end{align*}
Thus, $|f(x) - f(2)| < \epsilon$. Therefore, by the definition of continuity, $f(x)$ is continuous at $x=2$.
\end{proof}
\newpage
\section{Direct Proof}
\begin{center}(Problem $\#5$ Homework $\#3$)\end{center}
\noindent\textbf{Theorem:} For all positive real numbers $x$, the sum of $x$ and its reciprocal is greater than or equal to $2$.

\begin{proof}
Let $x \in \mathbb{R}^+$. Since $x$ is a real number, $(x-1)$ is also a real number, by closure under the real numbers.  A fundamental property of the real numbers is that the square of any real number is non-negative. Therefore, we establish that:
\[
    (x-1)^2 \ge 0
\]
Expanding the left-hand side yields:
\[
    x^2 - 2x + 1 \ge 0
\]
By adding $2x$ to both sides of the inequality, we obtain:
\[
    x^2 + 1 \ge 2x
\]
Because $x \in \mathbb{R}^+$, we know that $x$ is strictly positive ($x > 0$). Therefore, we can divide both sides of the inequality by $x$ without encountering zero-division or reversing the direction of the inequality sign. Doing so results in:
\begin{align*}
    \frac{x^2 + 1}{x} &\ge \frac{2x}{x} \\
    \frac{x^2}{x} + \frac{1}{x} &\ge 2 \\
    x + \frac{1}{x} &\ge 2
\end{align*}
Thus, for any positive real number $x$, the sum of $x$ and its reciprocal is greater than or equal to $2$.
\end{proof}

\end{document}
